% ============================================================================
\clearpage
\section*{v4.10.0-r30 Source-Only: One Number and Its Translations; Orientation of the Octant Step}
\addcontentsline{toc}{section}{v4.10.0-r30 Source-Only: One Number and Its Translations; Orientation of the Octant Step}
\label{sec:v4100-r30-translations}
% ============================================================================

\paragraph{Purpose.}
This layer collects the work after v4.10.0-r29, from 17 to 22 November 2026.  It has two parts.  The first part continues OP1: whether a boost can
change the area quantum, why frame T counts only \(j=\tfrac12\), the premise shared by all reasons for the assignment of the counting rules, a chance
estimate of the hits of \(\gamma_L\), a map of how the one number of OP1 enters OP2--OP5, and a profile of a possible mediator.  The second part
records the pre-registered tests H113--H117 on the bent octant of OP3, which locate the open choice in the orientation of the ordered return.  A short
remark lists the OP5 explorations and the computational checks of the same period.  Every OP1 item is an exploration; H113--H117 were pre-registered
but are marked \emph{not blind}, because their outcome was analytically fixed before the computation.  Earlier layers are not rewritten.  The notes
are listed in \S\ref{subsec:v4100-r30-notes}.

\noindent\textbf{Source-only status.}
This v4.10.0-r30 layer is source-only.  It does not build a PDF, does not assign a new DOI,
and does not authorize a public release.  The v5.2.0 Version DOI
\href{https://doi.org/10.5281/zenodo.24294648}{10.5281/zenodo.24294648} and the Concept DOI
\href{https://doi.org/10.5281/zenodo.19210728}{10.5281/zenodo.19210728} are unchanged by it.

\subsection*{v4.10.0-r30 results, part I: the one number of OP1}
\label{subsec:v4100-r30-results-op1}

\paragraph{Remark v4.10.0-r30.1 (a unitary boost does not change the area quantum; literature).}
Rovelli and Speziale (gr-qc/0205108) show that a boosted observer measures a different area operator \(A'\), which does not commute with \(A\); if the
boost acts unitarily, \(A\) and \(A'\) have the same spectrum and the same smallest area.  Only the distribution of eigenvalues, and hence the mean
area, changes.  In r10 the assignment of the counting rules (DL to M, ABCK to T) was shown to be equivalent to a frame-dependent area quantum
\(\gamma_M/\gamma_T=\gamma_L\).  Under reading (B) of r29.5, M and \(\bar M\) are related by a genuine boost and therefore see the same \(\gamma\);
the duality part \(J_B\) leaves the geometry fixed (r10).  Consequently no decomposition of the frame map into boost and duality yields a
frame-dependent area quantum while the boost acts unitarily.  The assignment is compatible with standard loop quantum gravity only if
(a) the map \(M\to T\) is not unitary (the \(J_B\) image with imaginary rapidity \(i\pi/2\) is the only candidate), or (b) \(\gamma_T\) is read as a
mean rather than an eigenvalue parameter; (b) is not a counting rule and again requires the pair \((M,\bar M)\) with \(\eta\) taken from \(\gamma_L\).

\paragraph{Remark v4.10.0-r30.2 (a boosted observer prefers but does not select extremal states; computation).}
For a puncture of spin \(k\) in the lowest block of the principal series of \(SL(2,\mathbb C)\) with \(\gamma=\gamma_M\), a boost along the horizon
normal with rapidity \(\eta_M=\operatorname{artanh}\beta_{\rm sym}=0.6175\) keeps the extremal states \(|k,\pm k\rangle\) with probability
\(0.82\) to \(0.74\) for \(k=\tfrac12,\dots,3\), the other states with \(0.73\) to \(0.49\).  The order is as expected; the criteria for a sharp
selection, fixed before the computation, are missed in both directions.  The rule ``every spin with weight two'' (U2) therefore does not follow from
reading (B) alone.  Transverse boosts distinguish nothing, so even the preference needs the boost along the normal as an additional assumption.

\paragraph{Remark v4.10.0-r30.3 (the rule of frame T has no data-free reason; exploration).}
Five reasons were collected for the rule that T counts only \(j=\tfrac12\): the tachyon floor, irreducibility, ``one crossing is one strand'', open
string endpoints, and the dimension ladder of r29.3.  The first two are already burdened (r10; the assumption (A) of 16 November), the next two are
explicit placeholders (CT-002, CT-004).  The dimension ladder T, M, P on line, plane and space holds exactly for units
\(\gamma_0\in[1,1.0757)\cdot\gamma_T\); its lower end is \(\gamma_T\) itself.  For T the ladder is therefore circular unless a unit in this window is
justified without reference to T.  Two candidates for such a reason were checked.  A square \(J_B^2=-1\) on extremal doublets selects all
half-integer \(j\) (Kramers, under (A)), every \(j\), or \(j=1,3,5,\dots\), depending on how \(J_B\) acts, but never \(j=\tfrac12\) alone.  A list
of 47 natural units, fixed before evaluation, has no entry in the window (0 found, about 0.94 expected).  The T rule remains a choice.

\paragraph{Remark v4.10.0-r30.4 (the premise shared by all reasons; exploration).}
Every reason for the assignment needs four parts: the area spectrum \(\propto\sqrt{j(j+1)}\) (K0, taken from the literature), weight two for M (K1),
all spins for M (K2), and \(j=\tfrac12\) only for T (K3).  Six routes were compared.  K0 is set in all of them, K1 has several partial supports, K2
is supported by none, and K3 only under (A) or the placeholders CT-002 and CT-004.  The weakest point is K3 (r30.3).

\paragraph{Remark v4.10.0-r30.5 (chance estimate of the hits of \(\gamma_L\); procedure fixed before the computation).}
Against a fixed catalogue of 16 targets, eight forms and a \(2\sigma\) window, the pair DL/ABCK gives three hits (\(\sin^2\theta_{12}\) and
\(\Omega_m\) by \(\beta^2\), \(\cot\theta_W\) by \(\gamma_L\)).  Among the 36 counting-rule pairs only this one has three hits.  The probability of at
least three hits for a random \(\gamma_L\) is \(0.018\) to \(0.023\); multiplied by the 36 pairs it is \(0.64\) to \(0.81\).  By the reading fixed
in advance this is compatible with chance once the selection of the pair is counted; only if the pair were fixed without looking at the data would
the probability be about \(2\%\).  Added afterwards: the two targets near \(0.30\) are strongly correlated, so two of the three hits are nearly one.

\paragraph{Remark v4.10.0-r30.6 (one number, five translations; exploration).}
The numerical values in OP2--OP5 all go back to the one number of OP1, \(\beta_{\rm sym}=0.5494\) (equivalently \(\eta\), \(\gamma_L\),
\(\gamma_M\)).  The translations of that number into the sectors do not: they are five separate identifications, four of them chosen after the
fact.
\begin{center}
\small
\begin{tabular}{@{}>{\raggedright\arraybackslash}p{0.17\textwidth}>{\raggedright\arraybackslash}p{0.33\textwidth}>{\raggedright\arraybackslash}p{0.40\textwidth}@{}}
\toprule
problem & translation (set) & own remainder, independent of the number\\
\midrule
OP2 & spectral offset \(=\beta_{\rm sym}\) & existence and selfadjointness of the operator\\
OP3 & residual bend \(\Phi_e=\pi/4-\gamma_M/2\) (H102) & origin of three generations, location of the bend\\
OP4 & \(|U_{e2}|^2=\beta^2\) (H9); mask parameter \(=\eta\) (H45/H46) & TM2, mask form, \(\delta_{\rm CP}\), CKM\\
OP5 & \(\Omega_M=\beta^2\) (H96) & mechanism selecting the ratio; absolute value of \(\Lambda\)\\
electroweak & \(\cot\theta_W=\gamma_L\) (H15) & no \(SU(2)\times U(1)\) in reading (B)\\
\bottomrule
\end{tabular}
\end{center}
C5 and C6/T2 do not use the number.  A common origin of the five translations would require a mediator.

\paragraph{Remark v4.10.0-r30.7 (profile of a mediator; exploration).}
A mediator would have to (A1) yield one dimensionless value fixing \(\eta\), as output, (A2) couple M and T through a state of its own, not through a
unitary frame change (r30.1), (A3) be stable, (A4) produce the five translations of r30.6 from one coupling, and (A5) carry a scale if it is to be
testable as a particle.  For (A3): by H74, equal spectra in M and T are unstable against any direct coupling \(\varepsilon\), and stability requires
\(|\omega_1^2-\omega_2^2|>2\varepsilon\).  Since \(\omega_M/\omega_T=\gamma_L\neq1\), a direct mediator is compatible with stability, which does not
fix \(\gamma_L\).  Of the three deferred candidates, the string bridge I1 touches two translations (the T count of OP1 and the residual bend of OP3),
under CT-002 and CT-004; the connector k185 takes \(\eta\) as input; the condensate I4 brings a scale but no translation.  None meets the profile.

\subsection*{v4.10.0-r30 results, part II: orientation of the octant step (H113--H117)}
\label{subsec:v4100-r30-results-op3}

\paragraph{Theoremlet v4.10.0-r30.8 (H113: no colour readout of the generation raster; not blind).}
In \(G\otimes(A_0\oplus A_c)\) with three colour-trivial generation directions, the nontrivial colour centre acts on \(G\otimes A_c\) as \(zI\),
\(z=(-1+i\sqrt3)/2\).  Every colour-equivariant linear readout \(T\) into a singlet satisfies \(zT=T\) there, hence \(T=0\), also on the three
diagonal places.  The colour readout of the literal \(e_4,e_5,e_6\) diagonal is refuted.  With a dual colour index the trace is the only invariant
readout (rank one).  The colour-neutral remainder column has rank three and carries, in the reading of H26, one weak doublet per row.

\paragraph{Theoremlet v4.10.0-r30.9 (H114: the photon point does not fix the Bargmann phase; not blind).}
At the photon eigenpoint of convention~F (\(b=0\), \(\rho_P=I_2/2\)) the inputs leave the Bargmann phase of the r14.5 Gram matrix free; both
\(\varphi=0\) and \(\varphi=\pi/4\) are allowed.  \emph{If} a readout identifies the phase with the radial Bures distance of r17.2,
\(d_B(r)=\tfrac12\arccos r\), and keeps the overlaps, then exactly
\[
\det G(r)=\frac{\sqrt{1+r}-1}{2},\qquad \lim_{r\downarrow0}\frac{\det G(r)}{r}=\frac14,
\]
so every \(r>0\) has complex rank three and \(r=0\) exactly rank two.  In the cap picture of r14.6 the same assumption reads
\(\Omega_{\rm eff}(r)=\arcsin r\).  The identification itself is not derived.

\paragraph{Theoremlet v4.10.0-r30.10 (H115: the direct compensator has rank zero on closed odd states; not blind).}
For the closed states of the generation coupling the embedding satisfies \(JS=-S\), hence \(CS=(S+JS)/2=0\).  The pre-registered expectation H115a,
that the existing compensator removes exactly one real direction and has rank five on these states, is refuted; the closure projector \(P_{C_7}\) is
a different operator.  A real rank-five halving does exist: the three
octant rays \(+x,+y,+z\) have the ordered overlap product \((1+i)/4\) (phase \(\pi/4\)) and a complex null line \(L=\mathbb Cu\),
\(u=(-i,1,-(1-i)/\sqrt2)\); the reflection \(J_n=I-2P_n\) along one real direction \(n\) of \(L\) gives \(C_n=I-P_n\) of real rank five, with
\(DC_n=D\) for the real octant readout \(D\).  The choice between \(n\) and \(in\) is not fixed by \(J\): the phase direction on \(L\) is
additional data.

\paragraph{Theoremlet v4.10.0-r30.11 (H116: magnitudes and endpoints do not fix the sign of the octant step; not blind).}
The loops \(x\to y\to z\to x\) and \(x\to z\to y\to x\) have the same start and end ray and three transition probabilities \(\tfrac12\); their ordered
overlap products are \((1+i)/4\) and its conjugate, i.e.\ steps of \(+\pi/4\) and \(-\pi/4\).  Two constant step sequences \(\pm1\) modulo eight visit
all eight phases in one cycle; a uniform histogram of accumulated phases therefore says nothing about the independence of single steps.  The
Vernam-type persistence heuristic of \S19 does not supply the sign.  What is missing is the selection of the ordered return among admissible
continuations; a first discriminating comparison repeats one visible initial state with different admissible histories.

\paragraph{Theoremlet v4.10.0-r30.12 (H117: two windows; not blind).}
Let the first step store the admissible orientations as a mask \(a=(a_+,a_-)\) inside one provisional state, without emitting a sign, and let the
second step test the same candidates with a mask \(b\); the survivors are \(s_d=a_db_d\).  Of the 16 mask pairs, three give \(+\pi/4\), three give
\(-\pi/4\), nine are blocked and one is ambiguous; a path dropped in the first window cannot return.  This is compatible with the two-step
persistence of \S19 and with the unique-successor condition of 19Q, provided the return event is triggered by the second, unique outcome.  The HC
repair rules contain no quantity that decides between the two signs in the second window.  The comparison was run in Wolfram Language and, for the
four additional cases, in MMA and both HDL cores of the tape machine.

\paragraph{Remark v4.10.0-r30.13 (OP5 explorations and computational checks of the period).}
The OP5 work of 18--22 November is recorded in the notes \texttt{op5\_\allowbreak bifocal\_\allowbreak response\_\allowbreak 2026\_\allowbreak 11\_\allowbreak 18}
and \texttt{op5\_\allowbreak pairing\_\allowbreak readout\_\allowbreak 2026\_\allowbreak 11\_\allowbreak 18}.  Among them: a closed form of the odd Dirac-shell plateau on
the sphere,
\[
B_{L,\ell}=\frac{(\ell-1)\ell(\ell+1)(\ell+2)}{64(2\ell+1)}\left[\frac{(\ell-1)!!}{\ell!!}\right]^2,\qquad L\geq\frac{\ell-1}{2},
\]
for completely occupied positive shells at fixed area and odd \(\ell\geq3\); an exclusion of equal radial probe signals for strictly convex bifocal
shapes; and a background model with paired vacuum share and unpaired dust that satisfies the limiting cases but does not derive the H96 ratio from
the existing terms.  The tape machine of \S19K--19L was compared in 209 bounded cases between MMA and two SystemVerilog controllers, with agreement.
None of this derives the native OP5 coupling; H96 remains open.

\subsection*{v4.10.0-r30 frames}
\label{subsec:v4100-r30-frames}

\begin{itemize}[leftmargin=2em]
\item R: rest frame of the field, \(\phi=1\), rapidity \(0\).
\item M: observer at rapidity \(\eta\), weight two; related to \(\bar M\) by a genuine, unitary boost, hence the same area spectrum (r30.1).
\item T: \(J_B\) image with shift \(i\pi/2\); \(j=\tfrac12\) only, without a data-free reason (r30.3); the only possible non-unitary step.
\item P: null field, \(\phi=0\), all \(2j+1\).
\end{itemize}
The octant tests of part II are made in the generation sphere and do not compare quantities across frames.

\subsection*{v4.10.0-r30 register entries}
\label{subsec:v4100-r30-register}

\begin{center}
\small
\begin{tabular}{@{}>{\raggedright\arraybackslash}p{0.34\textwidth}>{\raggedright\arraybackslash}p{0.58\textwidth}@{}}
\toprule
statement & status after v4.10.0-r30\\
\midrule
frame-dependent area quantum (r10) & requires a non-unitary frame map or a mean-value reading (r30.1)\\
rule of frame T (\(j=\tfrac12\)) & chosen; dimension ladder circular for T; \(J_B^2=-1\) and 47 units give no reason (r30.3)\\
hits of \(\gamma_L\) & compatible with chance once the pair selection is counted (r30.5)\\
translations of the one number & five, set separately (r30.6); no mediator meets the profile (r30.7)\\
direct compensator on closed odd states & rank zero (r30.10); H115a refuted, real rank-five halving only with an extra phase choice\\
sign of the octant step & not fixed by magnitudes, endpoints or \(J\); selection of the ordered return open (r30.11, r30.12)\\
H96, H102, H112 & open\\
\bottomrule
\end{tabular}
\end{center}

\subsection*{v4.10.0-r30 notes}
\label{subsec:v4100-r30-notes}

\texttt{op4\_\allowbreak status\_\allowbreak audit} (2026-11-17), \texttt{chance\_\allowbreak estimate}, \texttt{op1\_\allowbreak common\_\allowbreak premise},
\texttt{op1\_\allowbreak u2\_\allowbreak boost}, \texttt{op1\_\allowbreak t\_\allowbreak rule}, \texttt{op1\_\allowbreak area\_\allowbreak quantum\_\allowbreak boost},
\texttt{op\_\allowbreak root\_\allowbreak map}, \texttt{op\_\allowbreak mediator\_\allowbreak profile}, \texttt{op5\_\allowbreak bifocal\_\allowbreak response},
\texttt{op5\_\allowbreak pairing\_\allowbreak readout} (2026-11-18), each under \texttt{notes/\allowbreak research/}; the pre-registrations and results of
H113--H117 under \texttt{llm/} (2026-11-22).  The method notes of r30.2, r30.3 and r30.5 were committed before the computations.

\subsection*{v4.10.0-r30 ledger and guard update}
\label{subsec:v4100-r30-ledger-guard}

\[
R_{\rm chk}^{v4100\text{-}r30}=(1,1,0)
\]
for (status recorded; computations checked in the cited scripts; the counting rules per frame, the translations and the selection of the ordered
return not derived).

The layer spends no guard:
\[
R_{\rm glob}^{(v4100\text{-}r30,+)}=(0,1,1,1,1),
\qquad
\Gamma_{\rm guard}^{v4100\text{-}r30}=(1,1,1,1,1),
\qquad
\beta_4=1/256\ \text{unspent}.
\]
The layer does not imply exact finite \(C5\) zero closure, exact all-readout terminality, \(C6/T2\), a global photon-clock/RH theorem,
or global all-readout closure.  OP1--OP5, C5 and C6/T2 remain open.
